\subsection{A direct calculation at the exceptional prime three}
\label{sec:rh-dyadic-bernoulli-unit}

The exceptional coefficient prime $3$ in the ordinary RH construction
leads to an odd character of power-of-two conductor.  Its required
unramified cohomology vanishes for an elementary arithmetic reason.
The result below is uniform through the entire cyclotomic tower.

For a nontrivial primitive Dirichlet character $\chi$ of conductor $M$,
use the convention
\[
                 B_{1,\chi}=\frac1M\sum_{a=1}^{M}a\chi(a).
                                                               \tag{DB1}
\]
All reductions in characteristic $3$ are taken at an arbitrary prime
above $3$ in the character-value field.

\begin{theorem}[Dyadic Bernoulli values at three]
\label{thm:dyadic-bernoulli-three-unit}
If $\chi$ is odd and primitive of conductor $2^n$, then $n\ge2$
and $B_{1,\chi}$ is a unit at every prime above $3$.
Consequently, for every $n\ge2$,
\[
       3\nmid h^-(\mathbf Q(\mu_{2^n})),\qquad
       \operatorname{Cl}(\mathbf Q(\mu_{2^n}))[3^\infty]^- =0.
                                                               \tag{DB2}
\]
The superscript minus denotes the part on which complex conjugation
acts by $-1$; no assertion about the full real class number is needed.
\end{theorem}

\begin{proof}
The character of conductor $4$ gives $B_{1,\chi}=-1/2$.
Suppose $n\ge3$, put $M=2^n$ and $m=2^{n-2}$, and write
$\zeta=\chi(5)$.  The decomposition
\[
 (\mathbf Z/2^n\mathbf Z)^\times=\langle-1\rangle\times
                                      \langle5\rangle
\]
and primitivity show that $\zeta$ has order $m$.
Let $a_j$ be the least positive residue of $5^j$ modulo $M$.
Pairing $a_j$ with $M-a_j$ in (DB1), and using
$\sum_{j=0}^{m-1}\zeta^j=0$, gives
\[
 B_{1,\chi}=\frac2M\sum_{j=0}^{m-1}a_j\zeta^j
           =\sum_{j=0}^{m/2-1}s_j\zeta^j,
 \qquad s_j=\frac{2(a_j-a_{j+m/2})}{M}\in\{1,-1\}.
                                                               \tag{DB3}
\]
Indeed $5^{m/2}\equiv1+M/2\pmod M$, so the two residues
in the numerator differ by exactly $M/2$ or $-M/2$.
For example, the congruence follows by induction from
$v_2(5^{2^u}-1)=u+2$.

The reduction of $\zeta$ still has order $m$, because reduction
is injective on roots of unity of order prime to $3$.
When $n=3$, (DB3) is $1$ or $-1$.  When $n=4$, it is
a linear polynomial with coefficients $1$ or $-1$ evaluated
at a root of the irreducible polynomial $X^2+1$ over
$\mathbf F_3$, and hence is nonzero.

For $n\ge5$ put $h=2^{n-5}$.  In $\mathbf F_3[X]$ one has
\[
 \Phi_m(X)=X^{4h}+1
    =(X^{2h}+X^h-1)(X^{2h}-X^h-1).                            \tag{DB4}
\]
Both displayed factors are irreducible.  Every root of
$\Phi_m$ has order $m=8h$, and the degree of its minimal
polynomial is
\[
                 \operatorname{ord}_{2^{n-2}}(3)=2^{n-4}=2h.
                                                               \tag{DB5}
\]
For completeness, $v_2(3^{2^u}-1)=u+2$ for $u\ge1$;
this follows from $3^2-1=8$ by repeated squaring.
Together with $3\not\equiv1\pmod4$, it proves (DB5).

Write the polynomial in (DB3) as
$A_0+X^hA_1+X^{2h}A_2+X^{3h}A_3$, where each $A_i$
has degree less than $h$ and every coefficient is $1$ or $-1$.
For $s\in\{1,-1\}$, reduction modulo
$P_s=X^{2h}+sX^h-1$ gives
\[
           A_0+A_2-sA_3+X^h(A_1-sA_2-A_3).                    \tag{DB6}
\]
This cannot vanish.  Otherwise, at each of the $h$ coefficient
positions the four signs $a_0,a_1,a_2,a_3$ would satisfy
\[
           a_0+a_2-sa_3=0,\qquad a_1-sa_2-a_3=0
                                \quad\hbox{in }\mathbf F_3.
\]
Three signs have sum zero in $\mathbf F_3$ precisely when
all three are equal.  The first equality therefore gives
$a_2=-sa_3$, whereas the second gives $-sa_2=-a_3$.
These imply $a_3=-a_3$, a contradiction.
Neither irreducible factor in (DB4) divides the polynomial
in (DB3).  Its value at the reduction of $\zeta$ is nonzero.
This proves the unit assertion at every prime above $3$.

Set $K=\mathbf Q(\mu_{2^n})$ and $K^+=\mathbf Q(\mu_{2^n})^+$.
The relative class-number formula is
\[
 \frac{h_K}{h_{K^+}}=Q_Kw_K
       \prod_{\substack{\chi\in X(K)\\\chi(-1)=-1}}
                                    \left(-\frac{B_{1,\chi}}2\right),
 \quad Q_K=[\mathcal O_K^\times:
                 \mu(K)\mathcal O_{K^+}^\times]\in\{1,2\}.
                                                               \tag{DB7}
\]
Here $X(K)$ uses each character with its primitive conductor;
the formula and the unit-index statement are
\cite[Theorem 4.2.20 and Lemma 4.2.16]{dyadic-bernoulli-sharifi}.
Every such odd conductor is a power of $2$, and $w_K=2^n$.
Every factor in (DB7) is a $3$-adic unit, so its left side is
a $3$-adic unit.

Extension and norm of ideals identify
$\operatorname{Cl}(K^+)[3^\infty]$ with
$\operatorname{Cl}(K)[3^\infty]^+$: their composite on the
former is multiplication by $2$, and their other composite
is $1+j$, where $j$ is complex conjugation.  Multiplication
by $2$ is invertible on a $3$-group.  The ratio of the orders
of the two full $3$-primary class groups is therefore the
order of the minus part.  Since that ratio is one, (DB2) follows.
\end{proof}

\begin{corollary}[The actual unramified cocycle vanishing]
\label{cor:dyadic-three-unramified-vanishing}
Let $k$ be a field of characteristic $3$ with the discrete topology.
Let $\gamma:G_{\mathbf Q}\longrightarrow k^\times$ be an odd
character that factors through $\operatorname{Gal}(K/\mathbf Q)$,
where $K=\mathbf Q(\mu_{2^n})$ for some $n\ge2$.
Let $z:G_{\mathbf Q}\longrightarrow k$ be a locally constant
cocycle, so
$z(gh)=z(g)+\gamma(g)z(h)$.
Suppose that for every finite inertia subgroup $I$ there is
a subgroup $H\subset I\cap\ker\gamma$ of finite index prime
to $3$ on which $z$ vanishes.  Then there is $b\in k$ with
\[
                         z(g)=\gamma(g)b-b
                                      \quad(g\in G_{\mathbf Q}).
                                                               \tag{DB8}
\]
The same statement applies to the inverse of $\gamma$.
\end{corollary}

\begin{proof}
The restriction $u=z|_{G_K}$ is a continuous homomorphism to
the additive group of $k$.  It has finite image, since $G_K$
is compact and $k$ is discrete.  For every inertia group
$I_K\subset G_K$, the index
$[I_K:I_K\cap H]$ divides $[I\cap\ker\gamma:H]$.
Here $I_K=I\cap G_K$ is normal in $I\cap\ker\gamma$, so
the asserted index divisibility also holds when $H$ is not normal.
The image $u(I_K)$ is a group of exponent $3$ whose order
divides that prime-to-$3$ index, and hence it is zero.
The maximal unramified abelian extension and the Artin map
therefore identify $u$ with a homomorphism from
$\operatorname{Cl}(K)[3^\infty]$ to $k$
\cite[Chapter V, Example 3.9]{dyadic-bernoulli-milne}.
There are no real places of $K$ requiring a further condition.

The cocycle identity gives
$u(jgj^{-1})=\gamma(j)u(g)=-u(g)$.
Artin reciprocity is compatible with conjugation.
By (DB2) the $3$-primary class group has only its plus part,
so this odd homomorphism is zero.  Thus $z$ descends to a
cocycle of the finite group $G=\operatorname{Gal}(K/\mathbf Q)$.
Its order is a power of $2$.  Averaging proves (DB8) explicitly:
take $b=-|G|^{-1}\sum_{g\in G}z(g)$ and sum the cocycle
identity over its second argument.  Replacing $\gamma$
by its inverse preserves oddness and all the hypotheses.
\end{proof}

\paragraph{Application to the ordinary RH character.}
Write the actual finite character as $\psi=\psi_2\psi_{\rm odd}$.
In the ordinary RH host its odd-conductor factors are quadratic,
so $\psi_{\rm odd}^2=1$.  At coefficient prime $3$ the square
of the residual cyclotomic character is also trivial.
Consequently the residual character in the original weight-$-2$
vanishing requirement is exactly
\[
              \gamma=\overline{\eta_4\psi_2^{-2}},
              \qquad \xi=\eta_4\psi_2^2,                    \tag{DB9}
\]
where $\eta_4$ is the nontrivial character modulo $4$ and
$\gamma=\bar\xi^{-1}$.  If $\psi_2$ has conductor $2^c$
with $c\ge4$, then $\xi$ is odd and primitive of conductor
$2^{c-1}$: squaring halves the order on $\langle5\rangle$,
and multiplication by $\eta_4$ changes only the value at $-1$.
The original large-conductor case has $c\ge14$.

Thus Theorem~\ref{thm:dyadic-bernoulli-three-unit} supplies the
original Bernoulli-unit statement at $\ell=3$, odd conductor
$f=1$, and the original residual character.  Its source threshold
would only require $64\le2^{c-1}$; the proof works even without
that threshold.  Corollary~\ref{cor:dyadic-three-unramified-vanishing}
supplies the associated unramified-cocycle assertion with the
original prime-to-$3$ inertia indices.  If the source specifies
an arbitrary subgroup of $G_{\mathbf Q}$, intersect it with
$I\cap\ker\gamma$; its relative index is precisely the index
used above.  The source cocycle is locally constant, not an
arbitrary abstract cocycle\cite{dyadic-bernoulli-source}.

This discharges the residual unramified-vanishing input in this
ordinary RH branch.  The integral Selmer length estimate still
requires its actual deformation and spectral comparison.
For a general Dirichlet target, the odd-conductor factor of
$\psi^2$ need not be trivial; (DB9) and this application must
then be checked again.  No twisted Bernoulli-unit assertion
for a nontrivial odd-conductor character is claimed here.
